\documentclass[aspectratio=169,10pt]{beamer}

% Compile with XeLaTeX or LuaLaTeX for best Vietnamese support:
%   xelatex group9_hust_beamer.tex
%   xelatex group9_hust_beamer.tex

\usepackage{iftex}
\ifPDFTeX
  \usepackage[utf8]{inputenc}
  \usepackage[T5]{fontenc}
  \usepackage[vietnamese]{babel}
\else
  \usepackage{fontspec}
  \usepackage[vietnamese]{babel}
  \IfFontExistsTF{Arial}{\setsansfont{Arial}}{\setsansfont{TeX Gyre Heros}}
  \renewcommand{\familydefault}{\sfdefault}
\fi

\usepackage{amsmath}
\usepackage{booktabs}
\usepackage{graphicx}
\usepackage{tikz}
\usepackage{xcolor}
\usepackage{array}
\usetikzlibrary{arrows.meta,positioning,fit}

% ---------------------------------------------------------------------------
% HUST-inspired visual style
% ---------------------------------------------------------------------------
\definecolor{HUSTRed}{HTML}{B01116}
\definecolor{HUSTDark}{HTML}{1F2933}
\definecolor{HUSTGold}{HTML}{F2A900}
\definecolor{HUSTGray}{HTML}{F5F6F8}
\definecolor{HUSTLine}{HTML}{D9DEE7}

\setbeamertemplate{navigation symbols}{}
\setbeamertemplate{headline}{}
\setbeamertemplate{footline}{
  \leavevmode%
  \hbox{%
    \begin{beamercolorbox}[wd=.72\paperwidth,ht=2.6ex,dp=1.2ex,leftskip=1.8em]{author in head/foot}%
      \usebeamerfont{author in head/foot}Hybrid Fashion Recommender System -- Group 9
    \end{beamercolorbox}%
    \begin{beamercolorbox}[wd=.28\paperwidth,ht=2.6ex,dp=1.2ex,rightskip=1.8em plus1fil]{date in head/foot}%
      \usebeamerfont{date in head/foot}\insertframenumber/\inserttotalframenumber
    \end{beamercolorbox}%
  }%
}

\setbeamercolor{normal text}{fg=HUSTDark,bg=white}
\setbeamercolor{frametitle}{fg=white,bg=HUSTRed}
\setbeamercolor{title}{fg=white,bg=HUSTRed}
\setbeamercolor{subtitle}{fg=white}
\setbeamercolor{author in head/foot}{fg=white,bg=HUSTDark}
\setbeamercolor{date in head/foot}{fg=white,bg=HUSTRed}
\setbeamercolor{block title}{fg=white,bg=HUSTRed}
\setbeamercolor{block body}{fg=HUSTDark,bg=HUSTGray}
\setbeamercolor{alerted text}{fg=HUSTRed}
\setbeamercolor{itemize item}{fg=HUSTRed}
\setbeamercolor{itemize subitem}{fg=HUSTGold}
\setbeamercolor{enumerate item}{fg=HUSTRed}
\setbeamerfont{title}{series=\bfseries,size=\Large}
\setbeamerfont{subtitle}{size=\normalsize}
\setbeamerfont{frametitle}{series=\bfseries,size=\large}
\setbeamerfont{block title}{series=\bfseries}

\setbeamertemplate{frametitle}{
  \nointerlineskip
  \begin{beamercolorbox}[wd=\paperwidth,ht=3.2ex,dp=1.1ex,leftskip=1.1em]{frametitle}
    \insertframetitle
  \end{beamercolorbox}
}
\setbeamertemplate{itemize item}{\color{HUSTRed}\small$\blacktriangleright$}
\setbeamertemplate{itemize subitem}{\color{HUSTGold}\scriptsize$\bullet$}

\newcommand{\hustbar}{\tikz[baseline=-0.5ex]\fill[HUSTGold] (0,0) rectangle (1.1cm,0.08cm);}
\newcommand{\placeholder}[2]{%
  \begin{tikzpicture}
    \node[
      draw=HUSTLine,
      fill=HUSTGray,
      minimum width=#1,
      minimum height=#2,
      align=center,
      rounded corners=2pt,
      inner sep=6pt
    ] {\small\textcolor{HUSTDark}{Placeholder}\\[-2pt]\scriptsize chèn ảnh / biểu đồ sau};
  \end{tikzpicture}
}
\newcommand{\metric}[3]{%
  \begin{block}{#1}
    \centering
    {\Large\bfseries #2}\\[-2pt]
    {\scriptsize #3}
  \end{block}
}

\title[Hybrid Fashion Recommender]{Hybrid Fashion Recommender System}
\subtitle{Hệ thống gợi ý thời trang đa phương thức cho H\&M}
\author[Group 9]{%
  Nguyễn Thế Khải -- 202400050\\
  Phạm Gia Linh -- 202416262\\
  Nguyễn Đăng Long -- 202400057
}
\institute[HUST]{Đại học Bách khoa Hà Nội \\ Nhập môn Trí tuệ Nhân tạo -- CTTN KHMT K69}
\date{\today}

\begin{document}

% ---------------------------------------------------------------------------
\begin{frame}[plain]
  \begin{tikzpicture}[remember picture,overlay]
    \fill[HUSTRed] (current page.south west) rectangle (current page.north east);
    \fill[HUSTDark] (current page.south west) rectangle ([xshift=0.34\paperwidth]current page.north west);
    \fill[HUSTGold] ([xshift=0.34\paperwidth]current page.south west) rectangle ([xshift=0.355\paperwidth]current page.north west);
  \end{tikzpicture}
  \vspace{1.2cm}
  \begin{columns}[T,totalwidth=\textwidth]
    \begin{column}{0.31\textwidth}
      {\color{white}\bfseries\Large HUST}\\[4pt]
      {\color{white}\small School of Information and Communication Technology}\\[1.0cm]
      \placeholder{3.2cm}{3.2cm}\\[-2pt]
      {\color{white}\scriptsize Có thể thay bằng logo HUST}
    \end{column}
    \begin{column}{0.64\textwidth}
      {\color{white}\Huge\bfseries Hybrid Fashion Recommender System}\\[8pt]
      {\color{white}\Large Hệ thống gợi ý thời trang đa phương thức}\\[12pt]
      \hustbar\\[14pt]
      {\color{white}\normalsize Nhập môn Trí tuệ Nhân tạo -- Group 9}\\[8pt]
      {\color{white}\small Nguyễn Thế Khải \quad Phạm Gia Linh \quad Nguyễn Đăng Long}\\[18pt]
      {\color{white}\small Dataset: H\&M Personalized Fashion Recommendations}
    \end{column}
  \end{columns}
\end{frame}

% ---------------------------------------------------------------------------
\begin{frame}{Agenda}
  \begin{columns}[T,totalwidth=\textwidth]
    \begin{column}{0.48\textwidth}
      \begin{enumerate}
        \item Bài toán và mục tiêu
        \item Dữ liệu H\&M
        \item Pipeline xử lý
        \item Baseline, MF và NCF
      \end{enumerate}
    \end{column}
    \begin{column}{0.48\textwidth}
      \begin{enumerate}
        \setcounter{enumi}{4}
        \item Hybrid model
        \item Huấn luyện và đánh giá
        \item Demo Streamlit
        \item Kết luận và hướng phát triển
      \end{enumerate}
    \end{column}
  \end{columns}
\end{frame}

% ---------------------------------------------------------------------------
\section{Problem}
\begin{frame}{Đặt Vấn Đề}
  \begin{columns}[T,totalwidth=\textwidth]
    \begin{column}{0.56\textwidth}
      \begin{block}{Bối cảnh}
        Thương mại điện tử thời trang có số lượng sản phẩm lớn, thay đổi nhanh theo mùa vụ và xu hướng.
      \end{block}
      \begin{itemize}
        \item Người dùng khó tự tìm sản phẩm phù hợp trong kho lớn.
        \item Collaborative Filtering hiệu quả khi có lịch sử mua hàng, nhưng yếu với cold start.
        \item Thời trang phụ thuộc mạnh vào yếu tố thị giác: màu sắc, kiểu dáng, texture.
      \end{itemize}
    \end{column}
    \begin{column}{0.40\textwidth}
      \placeholder{\linewidth}{4.7cm}
      \vspace{2mm}
      {\scriptsize Gợi ý: chèn ảnh minh họa catalog / app thương mại điện tử.}
    \end{column}
  \end{columns}
\end{frame}

% ---------------------------------------------------------------------------
\begin{frame}{Mục Tiêu Dự Án}
  \begin{columns}[T,totalwidth=\textwidth]
    \begin{column}{0.50\textwidth}
      \begin{block}{Mục tiêu mô hình}
        Xây dựng mô hình lai kết hợp tín hiệu hành vi, metadata và ảnh sản phẩm.
      \end{block}
      \begin{itemize}
        \item Baseline: popularity.
        \item Collaborative Filtering: MF, NCF.
        \item Hybrid: NCF + visual features + metadata.
        \item Tối ưu và so sánh bằng MAP@12.
      \end{itemize}
    \end{column}
    \begin{column}{0.45\textwidth}
      \begin{block}{Mục tiêu hệ thống}
        Demo cục bộ bằng Streamlit cho recommendation top-12.
      \end{block}
      \begin{itemize}
        \item Chọn hoặc nhập User ID.
        \item Hiển thị sản phẩm đề xuất kèm ảnh.
        \item Có fallback sang popularity baseline khi thiếu checkpoint.
      \end{itemize}
    \end{column}
  \end{columns}
\end{frame}

% ---------------------------------------------------------------------------
\section{Data}
\begin{frame}{Dữ Liệu H\&M}
  \begin{columns}[T,totalwidth=\textwidth]
    \begin{column}{0.55\textwidth}
      \begin{table}
        \small
        \begin{tabular}{p{0.28\linewidth}p{0.62\linewidth}}
          \toprule
          \textbf{Nguồn} & \textbf{Vai trò} \\
          \midrule
          Transactions & Lịch sử mua hàng theo thời gian, tạo implicit feedback. \\
          Articles & Metadata sản phẩm: nhóm hàng, màu, department, garment group. \\
          Images & Ảnh sản phẩm dùng để trích xuất visual embedding. \\
          \bottomrule
        \end{tabular}
      \end{table}
      \begin{itemize}
        \item Encode \texttt{customer\_id}, \texttt{article\_id} sang index liên tục.
        \item Lưu dữ liệu đã xử lý ở định dạng Parquet để tăng tốc I/O.
      \end{itemize}
    \end{column}
    \begin{column}{0.40\textwidth}
      \placeholder{\linewidth}{4.8cm}
      \vspace{2mm}
      {\scriptsize Gợi ý: chèn sơ đồ 3 bảng dữ liệu hoặc vài ảnh sản phẩm mẫu.}
    \end{column}
  \end{columns}
\end{frame}

% ---------------------------------------------------------------------------
\section{Pipeline}
\begin{frame}{Pipeline Tổng Thể}
  \centering
  \begin{tikzpicture}[
    node distance=0.75cm and 0.75cm,
    box/.style={draw=HUSTLine, fill=HUSTGray, rounded corners=2pt, align=center, minimum width=2.35cm, minimum height=0.82cm, font=\small},
    redbox/.style={box, draw=HUSTRed, very thick},
    arr/.style={-{Latex[length=2mm]}, thick, color=HUSTDark}
  ]
    \node[box] (raw) {Raw Data\\transactions/articles/images};
    \node[box, right=of raw] (prep) {Preprocess\\Polars + encode ID};
    \node[box, right=of prep] (features) {Feature Store\\Parquet + NPZ};
    \node[redbox, below=of features] (model) {Models\\Popularity / MF / NCF / Hybrid};
    \node[box, left=of model] (eval) {Evaluation\\MAP@12, HR@12, NDCG@12};
    \node[box, right=of model] (demo) {Streamlit Demo\\Top-12 products};

    \draw[arr] (raw) -- (prep);
    \draw[arr] (prep) -- (features);
    \draw[arr] (features) -- (model);
    \draw[arr] (model) -- (eval);
    \draw[arr] (model) -- (demo);
  \end{tikzpicture}
  \vspace{5mm}
  \begin{block}{Code modules}
    \small \texttt{src/data\_processing} \quad
    \texttt{src/features} \quad
    \texttt{src/models} \quad
    \texttt{src/evaluation} \quad
    \texttt{app/app.py}
  \end{block}
\end{frame}

% ---------------------------------------------------------------------------
\begin{frame}{Tiền Xử Lý Và Feature Engineering}
  \begin{columns}[T,totalwidth=\textwidth]
    \begin{column}{0.49\textwidth}
      \begin{block}{Data preprocessing}
        \begin{itemize}
          \item Đọc dữ liệu bằng Polars.
          \item Join transactions với customers và articles.
          \item Lọc và chia train/test theo thời gian.
          \item Tạo index số nguyên cho user/item.
        \end{itemize}
      \end{block}
    \end{column}
    \begin{column}{0.49\textwidth}
      \begin{block}{Content features}
        \begin{itemize}
          \item Visual: ResNet-50 pretrained, vector 2048 chiều.
          \item Metadata: encode categorical attributes của article.
          \item Lookup thiếu feature: dùng vector 0 để pipeline vẫn chạy.
        \end{itemize}
      \end{block}
    \end{column}
  \end{columns}
  \vspace{2mm}
  \placeholder{0.94\linewidth}{2.0cm}
\end{frame}

% ---------------------------------------------------------------------------
\section{Models}
\begin{frame}{Các Mô Hình So Sánh}
  \begin{table}
    \small
    \begin{tabular}{p{0.20\linewidth}p{0.33\linewidth}p{0.33\linewidth}}
      \toprule
      \textbf{Phương pháp} & \textbf{Ưu điểm} & \textbf{Hạn chế} \\
      \midrule
      Popularity & Nhanh, ổn định, baseline rõ ràng & Cá nhân hóa thấp \\
      Matrix Factorization & Học latent user/item từ lịch sử mua hàng & Yếu khi user/item ít tương tác \\
      NCF & Biểu diễn phi tuyến qua GMF + MLP & Tốn tài nguyên, dễ overfit \\
      Hybrid & Kết hợp CF với metadata và ảnh & Thiết kế và training phức tạp hơn \\
      \bottomrule
    \end{tabular}
  \end{table}
  \begin{block}{Chiến lược}
    Bắt đầu từ baseline đơn giản, sau đó tăng dần độ phức tạp để đo cải thiện thực sự.
  \end{block}
\end{frame}

% ---------------------------------------------------------------------------
\begin{frame}{Matrix Factorization}
  \begin{columns}[T,totalwidth=\textwidth]
    \begin{column}{0.54\textwidth}
      \begin{block}{Ý tưởng}
        Mỗi user và item được ánh xạ vào không gian latent; điểm tương tác là dot product cộng bias.
      \end{block}
      \begin{equation*}
        \hat{r}_{u,i} = \mathbf{p}_u^\top \mathbf{q}_i + b_u + b_i
      \end{equation*}
      \begin{itemize}
        \item Implemented in \texttt{src/models/matrix\_factorization.py}.
        \item Huấn luyện trên implicit feedback với negative sampling.
      \end{itemize}
    \end{column}
    \begin{column}{0.40\textwidth}
      \placeholder{\linewidth}{4.5cm}
      \vspace{2mm}
      {\scriptsize Gợi ý: chèn sơ đồ user/item embedding.}
    \end{column}
  \end{columns}
\end{frame}

% ---------------------------------------------------------------------------
\begin{frame}{Neural Collaborative Filtering}
  \begin{columns}[T,totalwidth=\textwidth]
    \begin{column}{0.55\textwidth}
      \begin{block}{NeuMF architecture}
        NCF dùng hai nhánh song song để học tương tác user--item.
      \end{block}
      \begin{itemize}
        \item GMF: nhân từng phần tử giữa user embedding và item embedding.
        \item MLP: concat embedding, đưa qua nhiều dense layer.
        \item Output: concat hai nhánh và dự đoán logit tương tác.
      \end{itemize}
      \begin{equation*}
        \hat{y}_{ui} = f\left([\mathbf{p}_u \odot \mathbf{q}_i;\ \mathrm{MLP}(\mathbf{p}'_u,\mathbf{q}'_i)]\right)
      \end{equation*}
    \end{column}
    \begin{column}{0.40\textwidth}
      \placeholder{\linewidth}{4.6cm}
      \vspace{2mm}
      {\scriptsize Gợi ý: chèn diagram GMF + MLP.}
    \end{column}
  \end{columns}
\end{frame}

% ---------------------------------------------------------------------------
\begin{frame}{Hybrid Model: NCF + Visual + Metadata}
  \begin{columns}[T,totalwidth=\textwidth]
    \begin{column}{0.58\textwidth}
      \begin{block}{Fusion}
        Ghép latent vector từ NCF với vector ảnh 2048 chiều và metadata feature, sau đó đưa qua dense layers.
      \end{block}
      \begin{equation*}
        \hat{y}_{ui} =
        \mathrm{MLP}_{fusion}
        \left(
        [\mathbf{h}_{NCF}; \mathbf{v}_{image}; \mathbf{m}_{metadata}]
        \right)
      \end{equation*}
      \begin{itemize}
        \item Fusion layers: 256 $\rightarrow$ 64 $\rightarrow$ 1.
        \item BatchNorm, ReLU, Dropout để giảm overfitting.
        \item Score cuối cùng qua sigmoid khi inference.
      \end{itemize}
    \end{column}
    \begin{column}{0.38\textwidth}
      \placeholder{\linewidth}{5.0cm}
      \vspace{2mm}
      {\scriptsize Gợi ý: chèn architecture hybrid model.}
    \end{column}
  \end{columns}
\end{frame}

% ---------------------------------------------------------------------------
\section{Training}
\begin{frame}{Huấn Luyện}
  \begin{columns}[T,totalwidth=\textwidth]
    \begin{column}{0.49\textwidth}
      \begin{block}{Training setup}
        \begin{itemize}
          \item Loss: BCEWithLogitsLoss cho implicit feedback.
          \item Negative sampling: 4 negative / 1 positive.
          \item Optimizer: Adam, learning rate mặc định $10^{-3}$.
          \item Batch size mặc định: 4096.
        \end{itemize}
      \end{block}
    \end{column}
    \begin{column}{0.49\textwidth}
      \begin{block}{Ổn định training}
        \begin{itemize}
          \item Early stopping theo validation loss.
          \item Lưu \texttt{hybrid\_best.pt} và \texttt{hybrid\_last.pt}.
          \item Mixed precision khi có CUDA.
          \item Gradient clipping.
        \end{itemize}
      \end{block}
    \end{column}
  \end{columns}
\end{frame}

% ---------------------------------------------------------------------------
\begin{frame}{Metric Đánh Giá}
  \begin{columns}[T,totalwidth=\textwidth]
    \begin{column}{0.52\textwidth}
      \begin{block}{MAP@12}
        Metric chính theo bài toán H\&M: đo chất lượng danh sách top-12 có xét thứ tự.
      \end{block}
      \begin{equation*}
        \mathrm{MAP@12} =
        \frac{1}{|U|}\sum_{u \in U}
        \frac{1}{\min(m_u,12)}
        \sum_{k=1}^{12} P_u(k)\,\mathrm{rel}_u(k)
      \end{equation*}
    \end{column}
    \begin{column}{0.43\textwidth}
      \begin{block}{Metric bổ sung}
        \begin{itemize}
          \item Hit Rate@12: có ít nhất một hit trong top-12.
          \item NDCG@12: thưởng cho hit ở vị trí cao hơn.
        \end{itemize}
      \end{block}
      \placeholder{\linewidth}{2.0cm}
    \end{column}
  \end{columns}
\end{frame}

% ---------------------------------------------------------------------------
\section{Results}
\begin{frame}{Kết Quả Thử Nghiệm Hiện Có}
  \begin{columns}[T,totalwidth=\textwidth]
    \begin{column}{0.33\textwidth}
      \metric{MF}{0.2585}{MAP@12 smoke test}
    \end{column}
    \begin{column}{0.33\textwidth}
      \metric{NCF}{0.4727}{MAP@12 smoke test}
    \end{column}
    \begin{column}{0.33\textwidth}
      \metric{Popularity}{0.1674}{MAP@12 baseline run}
    \end{column}
  \end{columns}
  \vspace{4mm}
  \begin{columns}[T,totalwidth=\textwidth]
    \begin{column}{0.52\textwidth}
      \begin{itemize}
        \item Các số này lấy từ các file CSV trong \texttt{reports/}.
        \item NCF vượt MF và popularity trong smoke test hiện có.
        \item Khi có kết quả final, thay bảng và hình tại slide này.
      \end{itemize}
    \end{column}
    \begin{column}{0.43\textwidth}
      \placeholder{\linewidth}{3.0cm}
      \vspace{1mm}
      {\scriptsize Chèn \texttt{model\_metric\_comparison.png} hoặc biểu đồ final.}
    \end{column}
  \end{columns}
\end{frame}

% ---------------------------------------------------------------------------
\section{Demo}
\begin{frame}{Demo Streamlit}
  \begin{columns}[T,totalwidth=\textwidth]
    \begin{column}{0.45\textwidth}
      \begin{block}{Luồng sử dụng}
        \begin{enumerate}
          \item Chọn hoặc nhập \texttt{user\_id}.
          \item Chọn Hybrid Model hoặc Popularity Baseline.
          \item Hệ thống score candidate items.
          \item Hiển thị top-12 sản phẩm.
        \end{enumerate}
      \end{block}
      \begin{itemize}
        \item Module chính: \texttt{app/app.py}.
        \item Inference wrapper: \texttt{src/models/inference\_pipeline.py}.
      \end{itemize}
    \end{column}
    \begin{column}{0.50\textwidth}
      \placeholder{\linewidth}{5.0cm}
      \vspace{2mm}
      {\scriptsize Gợi ý: chèn screenshot Streamlit demo.}
    \end{column}
  \end{columns}
\end{frame}

% ---------------------------------------------------------------------------
\section{Risks}
\begin{frame}{Rủi Ro Và Cách Xử Lý}
  \begin{table}
    \small
    \begin{tabular}{p{0.28\linewidth}p{0.60\linewidth}}
      \toprule
      \textbf{Rủi ro} & \textbf{Hướng xử lý} \\
      \midrule
      Tài nguyên RAM/GPU hạn chế & Lọc theo thời gian, downsample, dùng Parquet, batch nhỏ, mixed precision. \\
      Overfitting & Dropout, L2 regularization, early stopping, bắt đầu từ model đơn giản. \\
      Dữ liệu theo mùa vụ/xu hướng & Chia train/test theo thời gian, so sánh với popularity gần thời điểm test. \\
      Thiếu visual feature & Fallback vector 0, vẫn dùng CF và metadata. \\
      \bottomrule
    \end{tabular}
  \end{table}
\end{frame}

% ---------------------------------------------------------------------------
\begin{frame}{Kết Luận}
  \begin{columns}[T,totalwidth=\textwidth]
    \begin{column}{0.55\textwidth}
      \begin{block}{Điểm đã đạt}
        \begin{itemize}
          \item Thiết kế pipeline recommendation modular.
          \item Có baseline, MF, NCF và Hybrid model.
          \item Có metric MAP@12, Hit Rate@12, NDCG@12.
          \item Có demo Streamlit cho top-12 recommendation.
        \end{itemize}
      \end{block}
    \end{column}
    \begin{column}{0.40\textwidth}
      \begin{block}{Hướng phát triển}
        \begin{itemize}
          \item Tuning hybrid model trên tập lớn hơn.
          \item Thêm time-aware recommendation.
          \item Cá nhân hóa theo style hoặc nhóm sản phẩm.
          \item Đánh giá ablation: CF only vs image vs metadata.
        \end{itemize}
      \end{block}
    \end{column}
  \end{columns}
\end{frame}

% ---------------------------------------------------------------------------
\begin{frame}[plain]
  \begin{tikzpicture}[remember picture,overlay]
    \fill[HUSTRed] (current page.south west) rectangle (current page.north east);
    \fill[HUSTDark] (current page.south west) rectangle ([yshift=0.23\paperheight]current page.north east);
    \fill[HUSTGold] (current page.south west) rectangle ([yshift=0.015\paperheight]current page.south east);
  \end{tikzpicture}
  \vfill
  \centering
  {\color{white}\Huge\bfseries Thank You}\\[10pt]
  {\color{white}\Large Q\&A}\\[18pt]
  {\color{white}\normalsize Hybrid Fashion Recommender System -- Group 9}
  \vfill
\end{frame}

\end{document}
